\documentclass[aspectratio=169,10pt]{beamer}

% Shared Beamer setup for Fall 2026 course slides.
\mode<presentation>{
  \usetheme{Madrid}
  \usecolortheme{default}
}

\definecolor{CalPolyGreen}{HTML}{154734}
\definecolor{CalPolyGold}{HTML}{BD8B13}
\definecolor{DeckInk}{HTML}{1F2933}
\definecolor{DeckMuted}{HTML}{5F6B7A}

\setbeamercolor{structure}{fg=CalPolyGreen}
\setbeamercolor{palette primary}{bg=CalPolyGreen,fg=white}
\setbeamercolor{palette secondary}{bg=DeckInk,fg=white}
\setbeamercolor{palette tertiary}{bg=CalPolyGold,fg=DeckInk}
\setbeamercolor{frametitle}{bg=CalPolyGreen,fg=white}
\setbeamercolor{title}{fg=CalPolyGreen}
\setbeamercolor{normal text}{fg=DeckInk,bg=white}
\setbeamercolor{block title}{bg=CalPolyGreen,fg=white}
\setbeamercolor{block body}{bg=black!3,fg=DeckInk}

\setbeamertemplate{navigation symbols}{}
\setbeamertemplate{footline}[frame number]
\setbeamertemplate{itemize item}{\color{CalPolyGold}$\blacktriangleright$}
\setbeamertemplate{itemize subitem}{\color{CalPolyGreen}$\bullet$}

\newcommand{\CourseNumber}{Course}
\newcommand{\CourseTitle}{Course Title}
\newcommand{\LectureNumber}{0}
\newcommand{\LectureTitle}{Lecture Title}
\newcommand{\LectureDate}{Date}
\newcommand{\CourseUnit}{Unit}

\newcommand{\coursenumber}[1]{\renewcommand{\CourseNumber}{#1}}
\newcommand{\coursetitle}[1]{\renewcommand{\CourseTitle}{#1}}
\newcommand{\lecturenumber}[1]{\renewcommand{\LectureNumber}{#1}}
\newcommand{\lecturetitle}[1]{\renewcommand{\LectureTitle}{#1}}
\newcommand{\lecturedate}[1]{\renewcommand{\LectureDate}{#1}}
\newcommand{\courseunit}[1]{\renewcommand{\CourseUnit}{#1}}

\author{Kirk Duran}
\institute{California Polytechnic State University, San Luis Obispo}
\date{\LectureDate}

\newcommand{\makedecktitle}{
  \begin{frame}[plain]
    \vspace{0.25in}
    {\usebeamerfont{title}\color{CalPolyGreen}\LectureTitle\par}
    \vspace{0.18in}
    {\large \CourseNumber{}: \CourseTitle\par}
    \vspace{0.08in}
    {\large Lecture \LectureNumber{} -- \LectureDate\par}
    \vfill
    {\small Kirk Duran \textbar{} Cal Poly, San Luis Obispo\par}
  \end{frame}
}

\newcommand{\placeholdernote}{
  \begin{block}{Placeholder}
    Replace these bullets with the final lecture material, examples, and in-class activities.
  \end{block}
}

\AtBeginSection[]{
  \begin{frame}{Roadmap}
    \tableofcontents[currentsection]
  \end{frame}
}


\pdfstringdefDisableCommands{\def\translate#1{#1}}

\graphicspath{{../assets/lecture-09/}}

% If an image file is not yet available, the slide displays a labeled
% placeholder using the intended filename. Place the matching file in
% ../assets/lecture-09/ to replace the placeholder automatically.
\newcommand{\lectureimage}[2][1.75in]{%
  \IfFileExists{../assets/lecture-09/#2}{%
    \includegraphics[width=\linewidth,height=#1,keepaspectratio]{#2}%
  }{%
    \fbox{%
      \parbox[c][#1][c]{0.94\linewidth}{%
        \centering
        \small\textbf{Image Placeholder}\\[0.08in]
        \scriptsize\texttt{\detokenize{#2}}
      }%
    }%
  }%
}

\newcommand{\lectureplaceholder}[2][1.75in]{%
  \fbox{%
    \parbox[c][#1][c]{0.94\linewidth}{%
      \centering
      \small\textbf{Image Placeholder}\\[0.08in]
      \scriptsize #2
    }%
  }%
}

\setbeamersize{text margin left=0.42in,text margin right=0.42in}

\coursenumber{CS 1001}
\coursetitle{Introduction to Computing}
\lecturenumber{9}
\lecturetitle{Functions II --- Return Values}
\lecturedate{Fall 2026}
\courseunit{1. Computing and Program Execution}

\begin{document}

\makedecktitle

\begin{frame}{Today}
  \begin{itemize}
    \item Treat a value-returning function call as an expression that itself produces a value.
    \item Trace argument evaluation and positional parameter binding explicitly.
    \item Model one function call using a temporary local execution environment.
    \item Define the semantics of \texttt{return}: evaluate a value, end the call, and resume the caller.
    \item Distinguish returning a value from printing a value.
    \item Trace repeated calls to the same function without conflating their local bindings.
  \end{itemize}

  \begin{keyidea}
    Lecture 8 treated a function as an abstraction. Today we open one call and trace exactly how its value is produced.
  \end{keyidea}
\end{frame}

\sectionframe{Part 1: A Function Call Is an Expression}

\begin{frame}{Expressions Produce Values}
  \begin{columns}[T,onlytextwidth]
    \begin{column}{0.52\textwidth}
      \begin{itemize}
        \item An expression is evaluated to obtain a value.
        \item Assignment then binds a name to that resulting value.
        \item This model does not change when the expression is a function call.
      \end{itemize}

      \begin{block}{Arithmetic expression}
        \texttt{answer = 3 + 4}
      \end{block}

      \begin{block}{Function-call expression}
        \texttt{answer = double(4)}
      \end{block}
    \end{column}

    \begin{column}{0.42\textwidth}
      % Intended visual: two expression trees, one arithmetic and one function call,
      % each reducing to a single value before assignment.
      \lectureimage[2.65in]{slide4.png}
    \end{column}
  \end{columns}
\end{frame}

\begin{frame}{The Returned Value Can Participate in a Larger Expression}
  \begin{columns}[T,onlytextwidth]
    \begin{column}{0.47\textwidth}
      \begin{block}{Assignment}
        \texttt{x = double(4)}\\[0.06in]
        \texttt{double(4) -> 8}\\
        \texttt{x -> 8}
      \end{block}

      \begin{block}{Arithmetic context}
        \texttt{y = 3 + double(4)}\\[0.06in]
        \texttt{3 + 8 -> 11}\\
        \texttt{y -> 11}
      \end{block}
    \end{column}

    \begin{column}{0.47\textwidth}
      \begin{block}{Comparison context}
        \texttt{double(4) > 5}\\[0.06in]
        \texttt{8 > 5 -> True}
      \end{block}

      \begin{alertblock}{Boundary for today}
        We will use ordinary expressions as arguments. Function calls nested inside other function calls are deferred to Lecture 10.
      \end{alertblock}
    \end{column}
  \end{columns}
\end{frame}

\sectionframe{Part 2: The Function-Call Evaluation Model}

\begin{frame}{A Call Temporarily Changes the Flow of Execution}
  \begin{columns}[T,onlytextwidth]
    \begin{column}{0.54\textwidth}
      \begin{itemize}
        \item Initial code normally executes sequentially.
        \item At a function call, evaluation transfers to the function body.
        \item The body executes in its own local execution environment.
        \item When the call returns, the suspended caller resumes.
      \end{itemize}

      \begin{block}{Control-flow shape}
        \centering
        caller $\rightarrow$ function body $\rightarrow$ caller resumes
      \end{block}
    \end{column}

    \begin{column}{0.40\textwidth}
      % Intended visual: a control-flow arrow leaving the caller, entering a function,
      % then returning to the exact call site.
      \lectureimage[2.55in]{slide7.png}
    \end{column}
  \end{columns}
\end{frame}

\begin{frame}[shrink=4]{The Call-Evaluation Procedure}
  \begin{enumerate}
    \item Encounter the function-call expression.
    \item Evaluate each argument expression to obtain argument values.
    \item Create a new local execution environment for this call.
    \item Bind each parameter name to its corresponding argument value.
    \item Execute statements in the function body sequentially.
    \item When \texttt{return expression} is encountered, evaluate that expression.
    \item End the current call and make the resulting value available to the caller.
    \item Resume evaluation at the location of the original call.
  \end{enumerate}

  \begin{keyidea}
    This procedure is the mental model we will repeatedly test with manual traces and the debugger.
  \end{keyidea}
\end{frame}

\begin{frame}{Trace One Call from Start to Finish}
  \begin{block}{Program}
    \texttt{def double(n: int) -> int:}\\
    \hspace*{0.22in}\texttt{result = n * 2}\\
    \hspace*{0.22in}\texttt{return result}\\[0.08in]
    \texttt{x = 4}\\
    \texttt{answer = double(x)}
  \end{block}

  \begin{block}{Before the call}
    \texttt{x -> 4}\qquad\texttt{answer: not yet bound}
  \end{block}

  \begin{activity}
    Predict the values of \texttt{n}, \texttt{result}, the returned value, and finally \texttt{answer}.
  \end{activity}
\end{frame}

\begin{frame}{The Argument Value Becomes the Parameter Binding}
  \begin{columns}[T,onlytextwidth]
    \begin{column}{0.46\textwidth}
      \begin{block}{Caller}
        \texttt{x -> 4}
      \end{block}

      \begin{block}{Argument expression}
        \texttt{x} $\longrightarrow$ \texttt{4}
      \end{block}
    \end{column}

    \begin{column}{0.46\textwidth}
      \begin{block}{Call to \texttt{double}}
        \texttt{n -> 4}
      \end{block}

      \begin{itemize}
        \item The name \texttt{x} is not renamed to \texttt{n}.
        \item The expression \texttt{x} is evaluated.
        \item Its resulting value is bound to \texttt{n}.
      \end{itemize}
    \end{column}
  \end{columns}

  \begin{keyidea}
    Data crosses the function boundary as a value; caller variable names do not become parameter names.
  \end{keyidea}
\end{frame}

\begin{frame}{The Function Body Then Executes Sequentially}
  \begin{block}{At function entry}
    \texttt{n -> 4}
  \end{block}

  \begin{block}{Execute the first statement}
    \texttt{result = n * 2}\\[0.06in]
    \texttt{n * 2} $\longrightarrow$ \texttt{4 * 2} $\longrightarrow$ \texttt{8}\\[0.06in]
    therefore \texttt{result -> 8}
  \end{block}

  \begin{itemize}
    \item Parameters already have bindings when the function body begins.
    \item Local variables appear only when their assignment statements execute.
  \end{itemize}
\end{frame}

\begin{frame}{Returning Reconnects the Function with Its Caller}
  \begin{columns}[T,onlytextwidth]
    \begin{column}{0.53\textwidth}
      \begin{block}{Inside \texttt{double}}
        \texttt{return result}\\[0.06in]
        \texttt{result -> 8}\\[0.06in]
        returned value: \texttt{8}
      \end{block}

      \begin{block}{Back in the caller}
        \texttt{answer = double(x)}\\
        conceptually becomes\\
        \texttt{answer = 8}
      \end{block}
    \end{column}

    \begin{column}{0.41\textwidth}
      % Intended visual: value 8 leaving a local function environment and flowing
      % back into the suspended assignment in the caller.
      \lectureimage[2.55in]{slide12.png}
    \end{column}
  \end{columns}
\end{frame}

\sectionframe{Part 3: Arguments and Parameters}

\begin{frame}{Arguments Are Expressions; Parameters Are Local Names}
  \begin{columns}[T,onlytextwidth]
    \begin{column}{0.46\textwidth}
      \begin{cpdefinition}
        An \textbf{argument} is an expression supplied at a function-call site.
      \end{cpdefinition}

      \begin{block}{Call site}
        \texttt{difference(x, y)}
      \end{block}
    \end{column}

    \begin{column}{0.46\textwidth}
      \begin{cpdefinition}
        A \textbf{parameter} is a local name declared in a function definition.
      \end{cpdefinition}

      \begin{block}{Definition}
        \texttt{def difference(first, second):}
      \end{block}
    \end{column}
  \end{columns}

  \begin{keyidea}
    Arguments are evaluated; parameters receive bindings to the resulting values.
  \end{keyidea}
\end{frame}

\begin{frame}{Multiple Positional Arguments Bind by Position}
  \begin{block}{Function and call}
    \texttt{def difference(first: int, second: int) -> int:}\\
    \hspace*{0.22in}\texttt{return first - second}\\[0.08in]
    \texttt{answer = difference(10, 4)}
  \end{block}

  \begin{center}
    \renewcommand{\arraystretch}{1.35}
    \begin{tabular}{c|c|c}
      \textbf{Position} & \textbf{Argument value} & \textbf{Parameter binding} \\\hline
      1 & \texttt{10} & \texttt{first -> 10} \\
      2 & \texttt{4} & \texttt{second -> 4}
    \end{tabular}
  \end{center}
\end{frame}

\begin{frame}{Changing Argument Order Changes the Bindings}
  \begin{columns}[T,onlytextwidth]
    \begin{column}{0.46\textwidth}
      \begin{block}{Call A}
        \texttt{difference(10, 4)}\\[0.08in]
        \texttt{first -> 10}\\
        \texttt{second -> 4}\\[0.08in]
        result: \texttt{6}
      \end{block}
    \end{column}

    \begin{column}{0.46\textwidth}
      \begin{block}{Call B}
        \texttt{difference(4, 10)}\\[0.08in]
        \texttt{first -> 4}\\
        \texttt{second -> 10}\\[0.08in]
        result: \texttt{-6}
      \end{block}
    \end{column}
  \end{columns}

  \begin{alertblock}{Positional semantics}
    Python does not infer intended roles from the values. For positional arguments, position determines the parameter binding.
  \end{alertblock}
\end{frame}

\begin{frame}{Caller Names Do Not Need to Match Parameter Names}
  \begin{block}{Program}
    \texttt{def rectangle\_area(width: int, height: int) -> int:}\\
    \hspace*{0.22in}\texttt{return width * height}\\[0.08in]
    \texttt{w = 5}\\
    \texttt{h = 3}\\
    \texttt{area = rectangle\_area(w, h)}
  \end{block}

  \begin{block}{Data movement}
    \centering
    \texttt{w -> 5 -> width}
    \qquad
    \texttt{h -> 3 -> height}
  \end{block}

  \begin{keyidea}
    Names are meaningful within their own context. The values, not the caller's variable names, are passed into the call.
  \end{keyidea}
\end{frame}

\begin{frame}{Trace an Argument Expression Pedantically}
  \begin{block}{Statement}
    \centering
    \Large\texttt{answer = square(x + 2)}
  \end{block}

  \begin{enumerate}
    \item Evaluate \texttt{x}: obtain \texttt{3}.
    \item Evaluate \texttt{3 + 2}: obtain \texttt{5}.
    \item Bind \texttt{n -> 5} for this call.
    \item Evaluate \texttt{n * n}: obtain \texttt{25}.
    \item Return \texttt{25} to the caller.
    \item Bind \texttt{answer -> 25}.
  \end{enumerate}
\end{frame}

\sectionframe{Part 4: The Semantics of \texttt{return}}

\begin{frame}{\texttt{return} Ends the Current Function Call}
  \begin{itemize}
    \item Once an executed \texttt{return} is reached, the current call is complete.
    \item Control transfers back to the caller.
    \item The caller receives the returned value as the value of the call expression.
  \end{itemize}

  \begin{block}{Conceptual rule}
    \centering
    \texttt{return expression}
    $\Longrightarrow$
    evaluate expression $\rightarrow$ end call $\rightarrow$ resume caller
  \end{block}

  \begin{alertblock}{Terminology}
    Avoid saying that \texttt{return} ``prints the answer'' or ``stores the answer.'' Neither description captures its semantics.
  \end{alertblock}
\end{frame}

\begin{frame}{The Caller Determines What Happens to the Returned Value}
  \begin{columns}[T,onlytextwidth]
    \begin{column}{0.30\textwidth}
      \begin{block}{Store it}
        \texttt{x = square(4)}\\[0.08in]
        \texttt{x -> 16}
      \end{block}
    \end{column}

    \begin{column}{0.30\textwidth}
      \begin{block}{Use it}
        \texttt{y = 2 * square(4)}\\[0.08in]
        \texttt{y -> 32}
      \end{block}
    \end{column}

    \begin{column}{0.30\textwidth}
      \begin{block}{Ignore it}
        \texttt{square(4)}\\[0.08in]
        value produced, but not bound
      \end{block}
    \end{column}
  \end{columns}

  \begin{keyidea}
    A function returns a value to its caller; the surrounding caller expression determines how that value is used.
  \end{keyidea}
\end{frame}

\begin{frame}{Printing and Returning Are Different Operations}
  \begin{columns}[T,onlytextwidth]
    \begin{column}{0.46\textwidth}
      \begin{block}{Return a value}
        \texttt{def square\_return(n):}\\
        \hspace*{0.18in}\texttt{return n * n}
      \end{block}
      \begin{itemize}
        \item Produces a value for the caller.
        \item The caller can store or combine that value.
      \end{itemize}
    \end{column}

    \begin{column}{0.46\textwidth}
      \begin{block}{Print a value}
        \texttt{def square\_print(n):}\\
        \hspace*{0.18in}\texttt{print(n * n)}
      \end{block}
      \begin{itemize}
        \item Produces visible output as a side effect.
        \item Does not explicitly return the printed number.
      \end{itemize}
    \end{column}
  \end{columns}
\end{frame}

\begin{frame}{A Function Without Explicit \texttt{return} Returns \texttt{None}}
  \begin{block}{Example}
    \texttt{def square\_print(n):}\\
    \hspace*{0.22in}\texttt{print(n * n)}\\[0.08in]
    \texttt{result = square\_print(4)}
  \end{block}

  \begin{columns}[T,onlytextwidth]
    \begin{column}{0.46\textwidth}
      \begin{block}{Visible output}
        \centering
        \texttt{16}
      \end{block}
    \end{column}

    \begin{column}{0.46\textwidth}
      \begin{block}{Returned value}
        \centering
        \texttt{None}
      \end{block}
    \end{column}
  \end{columns}

  \begin{keyidea}
    Visible output and a returned value are independent concepts.
  \end{keyidea}
\end{frame}

\begin{frame}{Prediction: Return or Print?}
  \begin{block}{Program}
    \texttt{def show\_double(n):}\\
    \hspace*{0.22in}\texttt{print(n * 2)}\\[0.08in]
    \texttt{def get\_double(n):}\\
    \hspace*{0.22in}\texttt{return n * 2}\\[0.08in]
    \texttt{a = show\_double(6)}\\
    \texttt{b = get\_double(6)}
  \end{block}

  \begin{activity}
    Predict (1) what is displayed, (2) the final binding of \texttt{a}, and (3) the final binding of \texttt{b}. Explain each using call semantics rather than intuition.
  \end{activity}
\end{frame}

\sectionframe{Part 5: Local Bindings and Independent Calls}

\begin{frame}{Each Call Has a Local Execution Environment}
  \begin{columns}[T,onlytextwidth]
    \begin{column}{0.52\textwidth}
      \begin{cpdefinition}
        A \textbf{local execution environment} is the set of parameter and local-variable bindings associated with one active function call.
      \end{cpdefinition}

      \begin{block}{During \texttt{rectangle\_area(5, 3)}}
        \texttt{width -> 5}\\
        \texttt{height -> 3}\\
        \texttt{area -> 15}
      \end{block}
    \end{column}

    \begin{column}{0.40\textwidth}
      % Intended visual: caller environment beside a temporary function environment.
      \lectureimage[2.55in]{slide26.png}
    \end{column}
  \end{columns}
\end{frame}

\begin{frame}{Local Variables Appear as Their Assignments Execute}
  \begin{block}{Function}
    \texttt{def rectangle\_area(width, height):}\\
    \hspace*{0.22in}\texttt{area = width * height}\\
    \hspace*{0.22in}\texttt{return area}
  \end{block}

  \begin{center}
    \renewcommand{\arraystretch}{1.35}
    \begin{tabular}{l|l}
      \textbf{Execution point} & \textbf{Local bindings} \\\hline
      Function entry & \texttt{width -> 5, height -> 3} \\
      After assignment & \texttt{width -> 5, height -> 3, area -> 15} \\
      After return & call no longer active
    \end{tabular}
  \end{center}
\end{frame}

\begin{frame}{Function-Local Bindings Have Call-Limited Lifetime}
  \begin{itemize}
    \item Parameter bindings are created for a specific call.
    \item Local-variable bindings are created as statements in that call execute.
    \item When the call completes, that local execution environment is no longer active.
    \item The returned value may survive because the caller can bind or use it.
  \end{itemize}

  \begin{block}{Before, during, after}
    \centering
    caller bindings
    $\rightarrow$
    caller + local call bindings
    $\rightarrow$
    caller bindings + returned result
  \end{block}

  \begin{alertblock}{Scope preview}
    Lecture 10 will examine what happens when more than one function call is active at the same time.
  \end{alertblock}
\end{frame}

\begin{frame}{Calling the Same Function Twice Creates Two Separate Calls}
  \begin{block}{Program}
    \texttt{def increase(value: int) -> int:}\\
    \hspace*{0.22in}\texttt{result = value + 1}\\
    \hspace*{0.22in}\texttt{return result}\\[0.08in]
    \texttt{first = increase(4)}\\
    \texttt{second = increase(10)}
  \end{block}

  \begin{columns}[T,onlytextwidth]
    \begin{column}{0.46\textwidth}
      \begin{block}{First call}
        \texttt{value -> 4}\\
        \texttt{result -> 5}\\
        returns \texttt{5}
      \end{block}
    \end{column}

    \begin{column}{0.46\textwidth}
      \begin{block}{Second call}
        \texttt{value -> 10}\\
        \texttt{result -> 11}\\
        returns \texttt{11}
      \end{block}
    \end{column}
  \end{columns}
\end{frame}

\begin{frame}{Trace Tables Need a Location Column Once Functions Appear}
  \begin{center}
    \renewcommand{\arraystretch}{1.25}
    \begin{tabular}{c|l|l|l}
      \textbf{Step} & \textbf{Location} & \textbf{Statement} & \textbf{Relevant bindings} \\\hline
      1 & caller & \texttt{x = 4} & \texttt{x -> 4} \\
      2 & \texttt{double} & function entry & \texttt{n -> 4} \\
      3 & \texttt{double} & \texttt{result = n * 2} & \texttt{result -> 8} \\
      4 & \texttt{double} & \texttt{return result} & returns \texttt{8} \\
      5 & caller & assignment completes & \texttt{answer -> 8}
    \end{tabular}
  \end{center}

  \begin{keyidea}
    A useful trace records both \emph{what executes} and \emph{where it executes}.
  \end{keyidea}
\end{frame}

\sectionframe{Part 6: Manual Tracing and Debugger Verification}

\begin{frame}{Manual Prediction Comes Before the Debugger}
  \begin{columns}[T,onlytextwidth]
    \begin{column}{0.52\textwidth}
      \begin{enumerate}
        \item Predict argument values.
        \item Predict parameter bindings.
        \item Predict each local assignment.
        \item Predict the returned value.
        \item Predict the caller's resulting state.
      \end{enumerate}

      \begin{keyidea}
        The debugger is evidence used to test a mental model, not a replacement for reasoning.
      \end{keyidea}
    \end{column}

    \begin{column}{0.40\textwidth}
      % Intended visual: Predict -> Trace -> Debug -> Compare loop.
      \lectureimage[2.55in]{slide32.png}
    \end{column}
  \end{columns}
\end{frame}

\begin{frame}{\texttt{Step Into} Exposes Function-Call Execution}
  \begin{itemize}
    \item \textbf{Step Over} executes a function call without tracing through its body line by line.
    \item \textbf{Step Into} transfers debugger focus into the called function.
    \item Inside the function, inspect parameter bindings and local variables as they appear.
    \item When the function returns, observe execution resume at the caller.
  \end{itemize}

  \begin{block}{Lab question}
    Does the debugger's Variables frame match the bindings predicted by your manual trace?
  \end{block}
\end{frame}

\begin{frame}{Worked Trace: Discounted Price}
  \begin{block}{Program}
    \texttt{def discounted\_price(price: float, rate: float) -> float:}\\
    \hspace*{0.22in}\texttt{discount = price * rate}\\
    \hspace*{0.22in}\texttt{final = price - discount}\\
    \hspace*{0.22in}\texttt{return final}\\[0.08in]
    \texttt{original = 80.0}\\
    \texttt{sale = discounted\_price(original, 0.25)}
  \end{block}

  \begin{activity}
    Trace the complete call. Record the argument values, parameter bindings, \texttt{discount}, \texttt{final}, returned value, and final binding of \texttt{sale}.
  \end{activity}
\end{frame}

\begin{frame}[shrink=8]{Worked Trace: Reason from Values, Not Similar Names}
  \begin{block}{Program}
    \texttt{def difference(first: int, second: int) -> int:}\\
    \hspace*{0.22in}\texttt{result = first - second}\\
    \hspace*{0.22in}\texttt{return result}\\[0.08in]
    \texttt{first = 20}\\
    \texttt{second = 7}\\
    \texttt{answer = difference(second, first)}
  \end{block}

  \begin{activity}
    Determine the parameter bindings without relying on the repeated names. Begin by evaluating the two argument expressions in their written order.
  \end{activity}

  \begin{alertblock}{Expected discipline}
    Caller names and parameter names may be identical, different, or misleading. Positional evaluation still determines the bindings.
  \end{alertblock}
\end{frame}

\begin{frame}{Common Failure Modes in Function Tracing}
  \begin{center}
    \renewcommand{\arraystretch}{1.28}
    \begin{tabular}{p{0.43\textwidth}|p{0.49\textwidth}}
      \textbf{Incorrect model} & \textbf{Correction} \\\hline
      ``The argument variable becomes the parameter.'' & The argument expression is evaluated; its value is bound to the parameter. \\
      ``\texttt{return} prints the result.'' & \texttt{return} produces the value of the call for the caller. \\
      ``A local variable exists as soon as the function starts.'' & A local binding appears when its assignment executes. \\
      ``The next call remembers the previous locals.'' & Each call creates a distinct local execution environment. \\
      ``Assignment happens before the call finishes.'' & The right-hand-side call must finish before assignment can bind its result.
    \end{tabular}
  \end{center}
\end{frame}

\begin{frame}{The Boundary of Today's Model}
  \begin{block}{Today: at most one active user-defined function call}
    \centering
    caller
    $\longrightarrow$
    function
    $\longrightarrow$
    caller resumes
  \end{block}

  \begin{itemize}
    \item We can trace calls with multiple parameters and local variables.
    \item We can evaluate ordinary expressions before and after a call.
    \item We can call the same function repeatedly, one call after another.
  \end{itemize}

  \begin{alertblock}{Next conceptual problem}
    What if a function calls another function \emph{before the first call has finished}?
  \end{alertblock}
\end{frame}

\begin{frame}{Takeaways}
  \begin{itemize}
    \item A value-returning function call is an expression and therefore evaluates to a value.
    \item Argument expressions are evaluated before their values are bound to parameters.
    \item Positional arguments correspond to parameters by position, not by caller variable name.
    \item One call has a temporary local execution environment containing parameter and local-variable bindings.
    \item \texttt{return} evaluates an expression, ends the current call, and supplies that value to the caller.
    \item Printing a value is not equivalent to returning it; functions without explicit \texttt{return} return \texttt{None}.
    \item Separate calls to the same function create separate local bindings.
  \end{itemize}

  \begin{keyidea}
    Mental model: \textbf{evaluate arguments $\rightarrow$ bind parameters $\rightarrow$ execute locally $\rightarrow$ return a value $\rightarrow$ resume the caller}.
  \end{keyidea}
\end{frame}

\end{document}
